\section*{Problemas}

\subsection*{Enunciados de los problemas}

% Recordar
\begin{problema}
	\label{prob:b94325d}
	Sea el espacio muestral $S$ y dos subconjuntos $A, B \subset S$. Enumera explícitamente los $2^2 = 4$ elementos disjuntos que conforman la partición fundamental $\particion(A,B)$ del espacio $S$.
\end{problema}

% Comprender
\begin{problema}
	\label{prob:320d62d}
	Extiende la construcción de una partición al caso de tres subconjuntos $A, B, C \subset S$. ¿Cuántos elementos tiene la partición $\particion(A,B,C)$ y cuáles son explícitamente estos subconjuntos elementales disjuntos? Explica por qué el número de elementos se duplica al agregar un tercer conjunto.
\end{problema}

% Aplicar
\begin{problema}
	\label{prob:18eada9}
	En un grupo de 50 estudiantes, 35 practican fútbol, 25 practican básquetbol y 15 practican ambos deportes. Usando las variables de la partición fundamental $\particion(F,B)$, ¿cuántos estudiantes no practican ningún deporte?
\end{problema}

% Analizar
\begin{problema}
	\label{prob:6d1bd44}
	Utilizando el principio de inclusión-exclusión o la teoría de particiones, demuestra la identidad de cardinalidad para la unión de tres conjuntos $A, B, C \subset S$:
	\begin{align}
		\card{A \cup B \cup C} = \card{A} + \card{B} + \card{C} - \card{A \cap B} - \card{A \cap C} - \card{B \cap C} + \card{A \cap B \cap C}.
	\end{align}
\end{problema}

% Evaluar
\begin{problema}
	\label{prob:15ead08}
	Un equipo de analistas en un ensayo clínico epidemiológico con un total de $\card{S} = 500$ pacientes evaluó la presencia de tres síntomas de una patología nueva ($S_1, S_2, S_3$). El informe preliminar presenta la siguiente estadística censal:
	\begin{itemize}
		\item $\card{S_1} = 300, \quad \card{S_2} = 250, \quad \card{S_3} = 200$
		\item $\card{S_1 \cap S_2} = 180, \quad \card{S_1 \cap S_3} = 150, \quad \card{S_2 \cap S_3} = 130$
		\item $\card{S_1 \cap S_2 \cap S_3} = 30$
	\end{itemize}
	Mediante el desglose algebraico en los 8 elementos de la partición $\particion(S_1,S_2,S_3)$, evalúa rigurosamente por qué el informe contiene un error lógico inevitable que imposibilita que estos datos correspondan a un censo real. Identifica específicamente cuál elemento de la partición arroja una cardinalidad matemáticamente absurda.
\end{problema}

% Crear
\begin{problema}
	\label{prob:95d5458}
	Un equipo de producto en una app de streaming musical clasifica a sus 800 usuarios activos según tres hábitos de consumo: escuchan \emph{podcasts} ($P$), crean listas de reproducción propias ($L$) y usan el modo sin conexión ($O$). Diseñe un ejemplo original con cardinalidades hipotéticas para $\card{P}$, $\card{L}$, $\card{O}$ y sus intersecciones (simples, dobles y triple) tales que los datos sean lógicamente consistentes (a diferencia del problema del ensayo clínico). Calcule, mediante el desglose en los 8 elementos de $\particion(P,L,O)$, cuántos usuarios no tienen ninguno de los tres hábitos.
\end{problema}

\subsection*{Soluciones detalladas}

% Recordar
\begin{solproblema}[prob:b94325d]
	Para dos subconjuntos $A, B \subset S$, cada elemento del espacio $S$ puede pertenecer o no pertenecer (complemento) a cada uno de los dos conjuntos. Esto genera $2 \times 2 = 4$ intersecciones elementales mutuamente disjuntas cuya unión totaliza al espacio $S$:
	\begin{itemize}
		\item $E_0 = A' \cap B'$ (elementos exteriores a ambos conjuntos)
		\item $E_1 = A' \cap B$ (elementos exclusivos de $B$)
		\item $E_2 = A \cap B'$ (elementos exclusivos de $A$)
		\item $E_3 = A \cap B$ (elementos en la intersección común de ambos)
	\end{itemize}
	Por construcción, se verifica que $E_i \cap E_j = \emptyset$ para todo $i \neq j$ y que $\bigcup_{k=0}^3 E_k = S$.
\end{solproblema}

% Comprender
\begin{solproblema}[prob:320d62d]
	Para tres subconjuntos $A, B, C \subset S$, al incluir la membresía dual o binaria respecto a un tercer conjunto $C$ o su complemento $C'$, el número de regiones se duplica respecto al caso de dos conjuntos: $2^3 = 8$ elementos elementales disjuntos.

	Estos 8 subconjuntos, ordenados sistemáticamente en código binario (donde pertenecer al conjunto activa un bit y pertenecer al complemento lo apaga), son:
	\begin{itemize}
		\item $E_0 = A' \cap B' \cap C'$ (ninguno de los tres conjuntos)
		\item $E_1 = A' \cap B' \cap C$ (sólamente en $C$)
		\item $E_2 = A' \cap B \cap C'$ (sólamente en $B$)
		\item $E_3 = A' \cap B \cap C$ (simultáneamente en $B$ y $C$, pero no en $A$)
		\item $E_4 = A \cap B' \cap C'$ (sólamente en $A$)
		\item $E_5 = A \cap B' \cap C$ (simultáneamente en $A$ y $C$, pero no en $B$)
		\item $E_6 = A \cap B \cap C'$ (simultáneamente en $A$ y $B$, pero no en $C$)
		\item $E_7 = A \cap B \cap C$ (en la intersección simultánea de los tres conjuntos)
	\end{itemize}
	El número se duplica porque cada nuevo conjunto añade un bit binario independiente (pertenece / no pertenece) a la codificación de cada región: pasar de 2 a 3 conjuntos multiplica por 2 el número de combinaciones posibles, de $2^2=4$ a $2^3=8$.
\end{solproblema}

% Aplicar
\begin{solproblema}[prob:18eada9]
	Sea $F$ la colección de estudiantes que practican fútbol y $B$ el conjunto que practica básquetbol. A partir de la partición del espacio muestral en cuatro regiones disjuntas, definimos las cardinalidades elementales:
	\begin{align}
		x_3 &= \card{F \cap B} = 15 \\
		x_2 &= \card{F \backslash B} = \card{F} - \card{F \cap B} = 35 - 15 = 20 \\
		x_1 &= \card{B \backslash F} = \card{B} - \card{F \cap B} = 25 - 15 = 10
	\end{align}
	El total de estudiantes que practican al menos un deporte es la unión $\card{F \cup B} = x_1 + x_2 + x_3 = 10 + 20 + 15 = 45$.

	Por complementariedad respecto a los 50 estudiantes del grupo:
	\begin{align}
		x_0 &= \card{(F \cup B)'} = \card{S} - \card{F \cup B} = 50 - 45 = 5.
	\end{align}
	Exactamente 5 estudiantes no practican ningún deporte.
\end{solproblema}

% Analizar
\begin{solproblema}[prob:6d1bd44]
	Consideremos primero el teorema de adición para dos conjuntos $A, B$, con las cardinalidades de la partición $\particion(A,B)$: $x_1 = \card{A' \cap B}$, $x_2 = \card{A \cap B'}$, $x_3 = \card{A \cap B}$. Dado que los conjuntos elementales de una partición son mutuamente disjuntos, la cardinalidad de cualquier unión es la suma directa de las cardinalidades de sus regiones componentes, lo cual da $\card{A} + \card{B} - \card{A \cap B} = \card{A \cup B}$.

	Para tres conjuntos, expresamos la unión agrupando dos términos y aplicando el teorema anterior:
	\begin{align}
		\card{A \cup B \cup C} &= \card{(A \cup B) \cup C} = \card{A \cup B} + \card{C} - \card{(A \cup B) \cap C}.
	\end{align}
	Por la ley distributiva de la intersección respecto a la unión, $(A \cup B) \cap C = (A \cap C) \cup (B \cap C)$. Aplicando nuevamente la fórmula de dos conjuntos sobre esta intersección distribuida:
	\begin{align}
		\card{(A \cup B) \cap C} &= \card{A \cap C} + \card{B \cap C} - \card{A \cap B \cap C}.
	\end{align}
	Sustituyendo esto y $\card{A \cup B} = \card{A} + \card{B} - \card{A \cap B}$ en la primera ecuación:
	\begin{align}
		\card{A \cup B \cup C} &= \card{A} + \card{B} + \card{C} - \card{A \cap B} - \card{A \cap C} - \card{B \cap C} + \card{A \cap B \cap C}.
	\end{align}
	Queda demostrada algebraicamente la fórmula del principio de inclusión-exclusión para orden tres.
\end{solproblema}

% Evaluar
\begin{solproblema}[prob:15ead08]
	Para verificar la coherencia lógica y matemática del informe, calculamos las variables $x_k = \card{E_k}$ ($k=0,1,\dots,7$) para la partición generada por los tres síntomas ($S_1, S_2, S_3$):
	\begin{align}
		x_7 &= \card{S_1 \cap S_2 \cap S_3} = 30 \\
		x_6 &= \card{S_1 \cap S_2 \cap S_3'} = \card{S_1 \cap S_2} - x_7 = 180 - 30 = 150 \\
		x_5 &= \card{S_1 \cap S_2' \cap S_3} = \card{S_1 \cap S_3} - x_7 = 150 - 30 = 120 \\
		x_3 &= \card{S_1' \cap S_2 \cap S_3} = \card{S_2 \cap S_3} - x_7 = 130 - 30 = 100
	\end{align}
	Deducimos ahora el número de pacientes con exclusivamente el primer síntoma ($x_4 = \card{S_1 \cap S_2' \cap S_3'}$):
	\begin{align}
		\card{S_1} = x_7 + x_6 + x_5 + x_4 \implies 300 = 30 + 150 + 120 + x_4 \implies x_4 = 0 \text{ (hasta aquí sin contradicción).}
	\end{align}
	Deducimos ahora el número de pacientes con exclusivamente el segundo síntoma ($x_2 = \card{S_1' \cap S_2 \cap S_3'}$):
	\begin{align}
		\card{S_2} = x_7 + x_6 + x_3 + x_2 \implies 250 = 30 + 150 + 100 + x_2 \implies x_2 = -30.
	\end{align}
	\textbf{Conclusión:} el cálculo arroja una cardinalidad negativa ($x_2 = -30$) para la región de pacientes que solo manifiestan el síntoma $S_2$. En la axiomática del conteo y la teoría de conjuntos, la función de cardinalidad $\card{\cdot} \ge 0$ no admite valores negativos. Esto demuestra sin lugar a dudas que las cifras publicadas en el informe preliminar son inconsistentes entre sí y reflejan un error en la recolección o tabulación censal de las intersecciones.
\end{solproblema}

% Crear
\begin{solproblema}[prob:95d5458]
	\textbf{Ejemplo:} sea $\card{S}=800$, con $\card{P}=350$ (podcasts), $\card{L}=400$ (listas propias), $\card{O}=300$ (modo sin conexión), $\card{P \cap L}=150$, $\card{P \cap O}=120$, $\card{L \cap O}=180$ y $\card{P \cap L \cap O}=60$.

	Calculamos las 8 regiones de $\particion(P,L,O)$:
	\begin{align}
		x_7 &= \card{P \cap L \cap O} = 60 \\
		x_6 &= \card{P \cap L \cap O'} = \card{P \cap L} - x_7 = 150-60 = 90 \\
		x_5 &= \card{P \cap L' \cap O} = \card{P \cap O} - x_7 = 120-60 = 60 \\
		x_3 &= \card{P' \cap L \cap O} = \card{L \cap O} - x_7 = 180-60 = 120 \\
		x_4 &= \card{P \cap L' \cap O'} = \card{P} - (x_7+x_6+x_5) = 350-(60+90+60) = 140 \\
		x_2 &= \card{P' \cap L \cap O'} = \card{L} - (x_7+x_6+x_3) = 400-(60+90+120) = 130 \\
		x_1 &= \card{P' \cap L' \cap O} = \card{O} - (x_7+x_5+x_3) = 300-(60+60+120) = 60
	\end{align}
	Todas las cardinalidades elementales son no negativas, por lo que los datos son lógicamente consistentes. El número de usuarios sin ninguno de los tres hábitos es:
	\begin{align}
		x_0 = \card{S} - \sum_{k=1}^7 x_k = 800 - (140+130+60+120+60+90+60) = 800 - 660 = 140.
	\end{align}
\end{solproblema}
